\chapter{AI Model Documentation / Training Details}\label{app:aimodel}

RecaLL does not train or fine-tune any model. It uses existing large language models through their
APIs (or locally through Ollama) and controls their behaviour with prompts, retrieved context and
output validation. This appendix is therefore a model card for how the models are \emph{used}.

\begin{table}[htbp]
\centering
\caption{Model card for RecaLL's AI features}\label{tab:modelcard}
\small
\begin{tabularx}{\textwidth}{@{}l X@{}}
\toprule
\textbf{Item} & \textbf{Description}\\
\midrule
Intended use & Help a student study their own documents: Cornell notes, Socratic questions and their
evaluation, flashcard and quiz generation, document summaries and concepts, explanations of selected
text, and cited answers in the Knowledge Hub.\\
Not intended for & Grading real examinations, medical/legal advice, or any use without the source
document available for checking.\\
Models & Chosen by the student. Configured defaults: Gemini (a Flash model), OpenAI
\texttt{gpt-4o-mini}, Groq \texttt{openai/gpt-oss-20b}, Ollama \texttt{gemma3:4b} (local). Any model the
provider offers can be selected. The team used Gemini \texttt{gemini-3.6-flash} and local models through
Ollama; the measurements in \Cref{ch:test} used \texttt{gemini-3.6-flash}.\\
Training data & None collected; no training performed. The provider's own model card applies to each
model.\\
Inputs & Topic or chunk text from the student's document (passed as untrusted reference data), the
student's question or answer, and a prompt: from the registry in \texttt{app/prompt\_manager.py} for
the study features, or built in the service code for the \texttt{/api/v1} features.\\
Outputs & JSON validated against Pydantic schemas (after an automatic repair step for malformed JSON)
or streamed Markdown text. Knowledge Hub citations are checked against the retrieved chunks.\\
Parameters & Structured generation uses temperature 0.1 by default and streaming 0.7; maximum output
tokens 4{,}096 and 8{,}192 respectively (\texttt{BaseLLMProvider}).\\
Reliability controls & Retries with exponential back-off and a Gemini rate limiter (study features),
a deterministic mock adapter when no key is available (\texttt{/api/v1}), schema validation and JSON
repair, citation validation, and a token-usage log per feature.\\
Evaluation & On the team's documents: 7 of 7 structured replies valid; grading followed the score bands in
4 of 4 cases; 5 of 5 answerable Knowledge Hub questions cited the relevant pages; an out-of-scope question
was answered as not covered (\Cref{ch:test}).\\
Usage notes & Outputs should be checked against the source, which every note, card and answer links
to; local models run on the student's own hardware, and cloud use needs an API key and an internet
connection.\\
Ethical notes & Data is sent only to the provider the student selects; keys are encrypted with
AES-256-GCM in the BYOK vault; no data is used to train models by RecaLL. Provider data-retention policies apply to
cloud requests.\\
\bottomrule
\end{tabularx}
\end{table}

\paragraph{AI tools used to build RecaLL.} The Antigravity agent generated most of the code from
requirements and designs the team had prepared beforehand, and was used to review the application and
find bugs; all output was reviewed and tested by the team. The automated tests were likewise largely
written with the agent and reviewed by the team. During development the team used the Gemini API
(\texttt{gemini-3.6-flash}) and local models through Ollama.

\paragraph{Prompt registry.} The prompts used by the application are listed in \Cref{tab:prompts}.
Each can be overridden by the user through the \texttt{/api/prompts} routes; placeholders are validated
before an override is saved, and every prompt can be reset to its default.

\begin{table}[htbp]
\centering
\caption{Prompt templates in \texttt{app/prompt\_manager.py}}\label{tab:prompts}
\small
\begin{tabularx}{\textwidth}{@{}l X@{}}
\toprule
\textbf{Key} & \textbf{Used for}\\
\midrule
\texttt{segment\_prompt} & Document processing: extract atomic concepts and topics from a section\\
\texttt{markdown\_cleanup\_prompt} & Document processing: repair OCR defects, broken tables and bare LaTeX\\
\texttt{handwritten\_outline\_reconcile\_prompt} & Document processing: derive chapter and section titles from OCR heading candidates\\
\texttt{cornell\_scaffold\_prompt} & Notes: structured Cornell notes with cue questions\\
\texttt{socratic\_question\_prompt} & Drills: two diagnostic questions (causal mechanism, counterfactual)\\
\texttt{socratic\_evaluation\_prompt} & Drills: short-answer grading, knowledge gaps and misconceptions\\
\texttt{flashcard\_prompt} & Flashcards for a topic\\
\texttt{selection\_flashcard\_prompt} & Flashcards from a highlighted PDF passage\\
\texttt{explain\_prompt} & Explain a highlighted passage\\
\texttt{summary\_prompt} & Markdown summary of a section\\
\texttt{chat\_prompt} & AI tutor sidebar in the topic reading view\\
\bottomrule
\end{tabularx}
\end{table}
